% Book 3, Year 1, chapter 7 -- modular programs, makefiles and tests.
% Provenance (never printed): sources/book3/BRIEFS.md, chapter 7.
% ledger: b1:make-origin, b1:gcc-mmd, b1:gcov, b1:ndebug
\chapter{Modular Programs, Makefiles and Tests}\label{ch:b1:c-modules-make-tests}

The town library's loan checker started as a hundred lines in one file. Three years later it is six hundred lines,
still in one file: dates, parsing, fines and reports mixed together. Every change recompiles everything; every
function can touch every \omterm{def:b1:c-functions-and-stack:lifetime}{global variable}; nobody knows which part a change may break, and nobody dares to touch
it. This chapter gives a C program the structure that makes change safe: modules with an \omterm{def:b1:c-modules-make-tests:interface}{interface}, a build that
recompiles only what a change affects, and tests that say, after each change, whether the program still does what
it did.

\begin{recall}
A module groups the functions that serve one purpose behind an \omterm{def:b1:c-modules-make-tests:interface}{interface}, the part its users rely on; the rest
can change. A specification states what a function must do; a test runs it on a chosen input and compares the result
with the expected one, and a set of tests is run after every change.
\end{recall}

\section{Headers and implementation files}\label{sec:b1:c-modules-make-tests:headers}

A C module is two files. The \code{.h} file declares what users may call; the \code{.c} file defines it. Users
write \texttt{\#include "ratio.h"}, which the \omterm{def:b1:compilation-pipeline:source-file}{preprocessor} replaces by the text of the header
(\cref{ch:b1:compilation-pipeline}); the \omterm{def:b1:compilation-pipeline:compiler}{compiler} then knows the names and types it needs to check their calls, and the
\omterm{def:b1:compilation-pipeline:compiler}{linker} later joins the calls to the definitions.

\begin{definition}[Header file, include guard]\label{def:b1:c-modules-make-tests:header}
A \emph{header file}\index{header file} is a file of declarations (types, function \omterm{def:b1:c-functions-and-stack:declaration}{prototypes}, constants) meant to be
included by other files. An \emph{include guard}\index{include guard} is the pair \texttt{\#ifndef NAME},
\texttt{\#define NAME} at its start and \texttt{\#endif} at its end, which makes a second inclusion of the header in the same
file expand to nothing.
\end{definition}

\begin{definition}[Interface, implementation]\label{def:b1:c-modules-make-tests:interface}
The \emph{interface}\index{interface (of a module)} of a module is what its users may rely on: in C, the declarations of
its header and the comments that state their contracts. Its \emph{implementation}\index{implementation (of a module)}
is everything else: the definitions in its \code{.c} file, which may change without any user changing.
\end{definition}

\omcode{code/bachelor-1/07-c-modules-make-tests/listings/ratio.h}{2}{17}{The \omterm{def:b1:c-modules-make-tests:interface}{interface} of a module of exact
fractions}\label{lst:b1:c-modules-make-tests:ratio-h}

\begin{definition}[Opaque type]\label{def:b1:c-modules-make-tests:opaque}
An \emph{opaque type}\index{opaque type} is a type that the \omterm{def:b1:c-modules-make-tests:interface}{interface} declares without defining it, such as
\code{struct ratio;} in a header whose fields appear only in the \omterm{def:b1:c-modules-make-tests:interface}{implementation} file. Users can hold and pass \omterm{def:b1:pointers-and-arrays:pointer}{pointers}
to it, but cannot read its fields, create one on the stack, or know its size: everything goes through the module's
functions.
\end{definition}

\omcode{code/bachelor-1/07-c-modules-make-tests/listings/ratio.c}{2}{40}{The \omterm{def:b1:c-modules-make-tests:interface}{implementation}: the fields, a private
helper, the constructor}\label{lst:b1:c-modules-make-tests:ratio-c}

\omcode{code/bachelor-1/07-c-modules-make-tests/listings/ratio_demo.c}{3}{18}{A user that knows only the
\omterm{def:b1:c-modules-make-tests:interface}{interface}}\label{lst:b1:c-modules-make-tests:demo}

\omoutput{code/bachelor-1/07-c-modules-make-tests/outputs/ratio_demo.txt}{Fractions kept in lowest terms; a zero
denominator refused.}

The function \code{gcd} is \code{static}: its name is visible only inside \code{ratio.c}
(\cref{ch:b1:c-functions-and-stack}), so another module may have its own \code{gcd} without a clash at link time, and
no user can come to depend on it. Because the type is opaque, the module can change its representation (keep fractions
unreduced, store them in \code{long long}, cache a floating value) and \code{ratio_demo.c} needs no change; it only needs
recompiling against the new \omterm{def:b1:compilation-pipeline:compiler}{object file}, if even that. The price is one heap block per fraction and a function call
per field read, which is why the C library's own \code{FILE} is opaque but small structures such as a point usually
are not.

\section{Separate compilation and make}\label{sec:b1:c-modules-make-tests:make}

\begin{definition}[Separate compilation]\label{def:b1:c-modules-make-tests:separate}
\emph{Separate compilation}\index{separate compilation} is the compiling of each \code{.c} file of a program on its own
into an \omterm{def:b1:compilation-pipeline:compiler}{object file} (\code{gcc -c ratio.c} writes \code{ratio.o}), the \omterm{def:b1:compilation-pipeline:compiler}{object files} being linked into the executable
in a last step. After a change, only the \omterm{def:b1:compilation-pipeline:compiler}{object files} whose sources (or included headers) changed must be compiled
again.
\end{definition}

A program's files form a graph: each \omterm{def:b1:compilation-pipeline:compiler}{object file} depends on its \code{.c} file and on the headers it includes; each
executable on its \omterm{def:b1:compilation-pipeline:compiler}{object files}. Deciding by hand what to recompile is error-prone; the program \code{make} does it from
a description of that graph.

\begin{history}[\code{make}]
\code{make} was written by Stuart Feldman at Bell Laboratories. His description of it, dated August 1978, shipped
with the Seventh Edition of Unix and starts from the problem of this chapter: in a large program, ``a change to a
declaration can obsolete a dozen other files''. Its rules, a \omterm{def:b1:c-modules-make-tests:makefile}{target}, the files it depends on and the commands that
make it, are still those of every \code{make} in use.
\end{history}

\begin{definition}[Makefile, target, prerequisite, recipe]\label{def:b1:c-modules-make-tests:makefile}
A \emph{makefile}\index{makefile} is a file of rules read by \code{make}. Each rule names a
\emph{target}\index{target (of a rule)}, usually a file to build; its \emph{prerequisites}\index{prerequisite}, the
files it is built from; and a \emph{recipe}\index{recipe}, the shell commands that build it, each line starting with a
tab character. \code{make} builds a target when it does not exist or when one of its prerequisites, after being brought
up to date itself, was modified more recently than the target.
\end{definition}

\omcode{code/bachelor-1/07-c-modules-make-tests/listings/Makefile}{2}{26}{A \omterm{def:b1:c-modules-make-tests:makefile}{makefile} for the module, a demo and
its tests}\label{lst:b1:c-modules-make-tests:makefile}

\code{CC} and \code{CFLAGS} are variables, used as \code{$(CC)}. The pattern rule \texttt{\%.o: \%.c} says how to make
any \omterm{def:b1:compilation-pipeline:compiler}{object file} from the \code{.c} file of the same name; in its \omterm{def:b1:c-modules-make-tests:makefile}{recipe} \code{$<} is the first \omterm{def:b1:c-modules-make-tests:makefile}{prerequisite} and, in the
\omterm{def:b1:compilation-pipeline:compiler}{linking} rules, \code{$@} is the \omterm{def:b1:c-modules-make-tests:makefile}{target} and \code{$^} all the \omterm{def:b1:c-modules-make-tests:makefile}{prerequisites}. The lines \code{ratio.o: ratio.h} add
\omterm{def:b1:c-modules-make-tests:makefile}{prerequisites} without a \omterm{def:b1:c-modules-make-tests:makefile}{recipe}: the headers each \omterm{def:b1:compilation-pipeline:compiler}{object file} includes. \code{all}, \code{test} and \code{clean} name no
file; \code{.PHONY} says so, so that \code{make} runs them even if a file of that name exists.

\omoutput{code/bachelor-1/07-c-modules-make-tests/outputs/make.txt}{Five commands, then none, then two, then all five:
\code{make} follows the timestamps.}\label{out:b1:c-modules-make-tests:make}

\begin{omfigure}
\begin{tikzpicture}[f/.style={draw, rounded corners=2pt, font=\small\ttfamily, minimum height=6mm, inner sep=3pt},
                    >={Stealth[length=2mm]}, x=2.6cm, y=1.25cm]
\node[f, fill=omCode!8] (h) at (1.6,1.0) {ratio.h};
\node[f] (dc) at (0,0) {ratio\_demo.c};
\node[f] (rc) at (1.38,0) {ratio.c};
\node[f, fill=omCode!8] (k) at (2.9,0) {check.h};
\node[f] (tc) at (3.9,0) {ratio\_test.c};
\node[f, fill=omDef!10] (do) at (0.3,-1.3) {ratio\_demo.o};
\node[f, fill=omDef!10] (ro) at (1.9,-1.3) {ratio.o};
\node[f, fill=omDef!10] (to) at (3.3,-1.3) {ratio\_test.o};
\node[f, fill=omThm!15] (de) at (1.0,-2.5) {ratio\_demo};
\node[f, fill=omThm!15] (te) at (2.7,-2.5) {ratio\_test};
\foreach \a/\b in {dc/do, h/do, rc/ro, h/ro, tc/to, h/to, k/to, do/de, ro/de, ro/te, to/te}
  {\draw[->] (\a) -- (\b);}
\node[font=\scriptsize, anchor=west] at (4.4,-1.3) {object files};
\node[font=\scriptsize, anchor=west] at (4.4,-2.5) {executables};
\node[font=\scriptsize, align=left, anchor=west] at (4.4,0.5) {sources and\\headers};
\end{tikzpicture}
\par\smallskip{\small The dependency graph that \cref{lst:b1:c-modules-make-tests:makefile} describes, arrows from
\omterm{def:b1:c-modules-make-tests:makefile}{prerequisite} to \omterm{def:b1:c-modules-make-tests:makefile}{target}. A change to a file must reach every file downstream of it, and only those.}
\end{omfigure}

\begin{omfigure}
\begin{tikzpicture}[f/.style={draw, rounded corners=2pt, font=\small\ttfamily, minimum height=6mm, inner sep=3pt,
                    align=center}, >={Stealth[length=2mm]}, x=3.3cm, y=1.45cm]
\node[f, fill=omThm!20] (dc) at (0,0) {ratio\_demo.c\\\scriptsize touched: newest};
\node[f] (h) at (1,0) {ratio.h};
\node[f] (rc) at (2,0) {ratio.c};
\node[f, fill=omThm!20] (do) at (0.5,-1.2) {ratio\_demo.o\\\scriptsize older than a prerequisite};
\node[f] (ro) at (1.9,-1.2) {ratio.o\\\scriptsize newer than both};
\node[f, fill=omThm!20] (de) at (1,-2.4) {ratio\_demo\\\scriptsize older than ratio\_demo.o};
\foreach \a/\b in {dc/do, h/do, rc/ro, h/ro, do/de, ro/de} {\draw[->] (\a) -- (\b);}
\end{tikzpicture}
\par\smallskip{\small \code{make} after \code{touch ratio_demo.c}. Visiting the graph from the \omterm{def:b1:c-modules-make-tests:makefile}{target} down, it finds
\code{ratio_demo.o} older than one of its \omterm{def:b1:c-modules-make-tests:makefile}{prerequisites} and rebuilds it; \code{ratio_demo}, now older than that new
\omterm{def:b1:compilation-pipeline:compiler}{object file}, is linked again; \code{ratio.o} is left alone. Shaded: the two commands of the third run.}
\end{omfigure}

The \omterm{def:b1:c-modules-make-tests:makefile}{prerequisites} \code{ratio.o: ratio.h} are written by hand, and a forgotten one is a classic bug: the \omterm{def:b1:compilation-pipeline:compiler}{object file} is
not rebuilt after a header changes, and the program silently runs with old values
(\cref{exo:b1:c-modules-make-tests:5}). The \omterm{def:b1:compilation-pipeline:compiler}{compiler} can write these lines itself: with \code{-MMD}, compiling
\code{loans.c} also writes \code{loans.d}, a \omterm{def:b1:c-modules-make-tests:makefile}{makefile} fragment listing the headers \code{loans.c} included, which the
\omterm{def:b1:c-modules-make-tests:makefile}{makefile} includes (\cref{exo:b1:c-modules-make-tests:11}).

\section{Unit tests in C}\label{sec:b1:c-modules-make-tests:unit-tests}

\begin{definition}[Unit test, test suite]\label{def:b1:c-modules-make-tests:unit-test}
A \emph{unit test}\index{unit test} checks one function or one module on its own, through its \omterm{def:b1:c-modules-make-tests:interface}{interface}, on chosen
inputs, against expected results. A \emph{test suite}\index{test suite} is the set of tests of a program, run as a
whole after every change; a single failure means the change is not ready.
\end{definition}

C has no test framework in its library. A \omterm{def:b1:compilation-pipeline:source-file}{macro} of a few lines is enough to start: it counts the checks, prints each
failure with its file, line and the text of the condition (\texttt{\#cond} turns the \omterm{def:b1:compilation-pipeline:source-file}{macro}'s argument into a string), and
goes on, so that one run reports every failure.

\omcode{code/bachelor-1/07-c-modules-make-tests/listings/check.h}{2}{20}{A check that counts and goes
on}\label{lst:b1:c-modules-make-tests:check}

\omcode{code/bachelor-1/07-c-modules-make-tests/listings/ratio_test.c}{4}{35}{The tests of the ratio
module}\label{lst:b1:c-modules-make-tests:ratio-test}

\omoutput{code/bachelor-1/07-c-modules-make-tests/outputs/ratio_test.txt}{Nine checks.}

The test program returns a nonzero exit status when a check fails, so that \code{make test} and any script that runs
it stop there. A failing test is the best bug report: it gives the input, the expected result and the place. When a bug
is found, write the test that shows it first, watch it fail, then fix the code and watch it pass: the test then guards
against the bug's return. Tests are compiled with \code{-Werror}, which turns every warning into an error.

\section{Designing tests}\label{sec:b1:c-modules-make-tests:design}

A \omterm{def:b1:c-modules-make-tests:unit-test}{test suite} cannot try every input: \code{classify} below has $2^{96}$ possible triples. The inputs are chosen so
that every behaviour of the function is tried at least once, and every place where its behaviour changes is tried on
both sides.

\begin{definition}[Input class, boundary value]\label{def:b1:c-modules-make-tests:input-class}
An \emph{input class}\index{input class} is a set of inputs that the specification treats the same way, so that one
test of the class stands for all of them; the classes of a function form a partition of its domain. A \emph{boundary
value}\index{boundary value} is an input at the edge of a class, next to an input of another class, where off-by-one
errors live.
\end{definition}

\omcode{code/bachelor-1/07-c-modules-make-tests/listings/classify.c}{2}{15}{The kind of a triangle with sides $a$, $b$
and $c$}\label{lst:b1:c-modules-make-tests:classify}

The classes of \code{classify} follow from its specification: a nonpositive side; positive sides that break the
triangle inequality (the longest side as $a$, as $b$, as $c$); an equilateral triangle; isosceles with the equal pair in
each of three places; scalene. The boundaries: a side equal to 0; the flat triple \code{1 2 3}, where $a + b = c$
exactly. The second kind of guidance comes from the code itself.

\begin{definition}[Control-flow graph, feasible path]\label{def:b1:c-modules-make-tests:cfg}
The \emph{control-flow graph}\index{control-flow graph} of a function has a node for each block of statements executed
in sequence, and an arc from a block to each block that can run next; a node with two outgoing arcs is a decision. A
\emph{path} runs from the entry to an exit; it is \emph{feasible}\index{feasible path} if some input makes the function
follow it.
\end{definition}

\begin{definition}[Statement, branch and path coverage]\label{def:b1:c-modules-make-tests:coverage}
A \omterm{def:b1:c-modules-make-tests:unit-test}{test suite} reaches \emph{statement coverage}\index{statement coverage} of a function if its tests together execute
every statement; \emph{branch coverage}\index{branch coverage} if they take every arc of every decision, both the
true and the false one; \emph{path coverage}\index{path coverage} if they follow every \omterm{def:b1:c-modules-make-tests:cfg}{feasible path}.
\end{definition}

\begin{omfigure}
\begin{tikzpicture}[d/.style={draw, diamond, aspect=2.6, font=\scriptsize, inner sep=1pt},
                    s/.style={draw, rounded corners=2pt, font=\scriptsize, inner sep=3pt, minimum height=5mm},
                    >={Stealth[length=1.8mm]}, y=1.15cm]
\node[d] (d1) at (0,0) {a,b,c $>$ 0?};
\node[d] (d2) at (0,-1.2) {inequality holds?};
\node[d] (d3) at (0,-2.4) {a=b=c?};
\node[d] (d4) at (0,-3.6) {two equal?};
\node[s] (s5) at (0,-4.8) {return SCALENE};
\node[s] (n1) at (4.2,0) {return NOT (line 7)};
\node[s] (n2) at (4.2,-1.2) {return NOT (line 9)};
\node[s] (e) at (4.2,-2.4) {return EQUILATERAL};
\node[s] (i) at (4.2,-3.6) {return ISOSCELES (line 13)};
\draw[->] (d1) -- node[left, font=\scriptsize] {yes} (d2);
\draw[->] (d2) -- node[left, font=\scriptsize] {yes} (d3);
\draw[->] (d3) -- node[left, font=\scriptsize] {no} (d4);
\draw[->] (d4) -- node[left, font=\scriptsize] {no} (s5);
\draw[->, dashed, omThm, thick] (d1) -- node[above, font=\scriptsize] {no} (n1);
\draw[->] (d2) -- node[above, font=\scriptsize] {no} (n2);
\draw[->] (d3) -- node[above, font=\scriptsize] {yes} (e);
\draw[->, dashed, omThm, thick] (d4) -- node[above, font=\scriptsize] {yes} (i);
\node[font=\scriptsize, align=left, anchor=west] at (-5.9,-1.2) {\code{3 3 3}: yes, yes, yes\\\code{3 4 5}: yes, yes,
no, no\\\code{1 2 8}: yes, no};
\end{tikzpicture}
\par\smallskip{\small The \omterm{def:b1:c-modules-make-tests:cfg}{control-flow graph} of \code{classify}, each compound condition drawn as one decision. The three
tests of the first run follow three paths; the dashed arcs are taken by none of them, so lines 7 and 13 never run.}
\end{omfigure}

\texttt{gcc -{}-{}coverage} adds counters to the compiled code; the program writes them to a file when it ends, and
\code{gcov} turns them into a copy of the source with each line's count, \texttt{\#\#\#\#\#} marking a line never run;
\code{-b} adds the branches.

\omoutput{code/bachelor-1/07-c-modules-make-tests/outputs/coverage.txt}{Three tests, then twelve: 80\,\% of the lines and
59\,\% of the branches, then all of them.}\label{out:b1:c-modules-make-tests:coverage}

\code{gcov} counts 22 branches for 4 decisions: it sees each atomic condition of \code{a <= 0 || b <= 0 || c <= 0} as
a decision of its own, because \code{||} evaluates its operands one at a time and stops at the first true one. Its
\omterm{def:b1:c-modules-make-tests:coverage}{branch coverage} is therefore finer than that of the graph above: to take all six arcs of line 6, each of the three
conditions must be true in some test where the earlier ones are false, which is what the tests \code{0 4 4},
\code{4 -1 4} and \code{4 4 0} of the second run do. Full coverage says that every part of the code ran, not that its
results were right: a coverage figure measures the tests, never the program.

\section{Assertions, defensive programming and the debugger}\label{sec:b1:c-modules-make-tests:assert}

\begin{definition}[Assertion]\label{def:b1:c-modules-make-tests:assertion}
An \emph{assertion}\index{assertion} is a condition that the programmer states must hold at a point of the program.
In C, \code{assert(cond)} from \code{<assert.h>} evaluates \code{cond}; if it is false, it prints the condition, the file,
the line and the function on standard error and stops the program with \code{abort}. If the \omterm{def:b1:compilation-pipeline:source-file}{macro} \code{NDEBUG} is
defined, every \code{assert} expands to nothing.
\end{definition}

\omcode{code/bachelor-1/07-c-modules-make-tests/listings/assert_demo.c}{2}{20}{A precondition stated as an
\omterm{def:b1:c-modules-make-tests:assertion}{assertion}}\label{lst:b1:c-modules-make-tests:assert}

\omoutput{code/bachelor-1/07-c-modules-make-tests/outputs/assert.txt}{With the check, a clear stop; without it,
\code{NaN} printed as a mean.}\label{out:b1:c-modules-make-tests:assert}

Exit status 134 is $128 + 6$, the number of the signal sent by \code{abort}. Built with \code{-DNDEBUG}, the program
divides 0 by 0 and prints \code{-nan}, ``not a number'', the floating-point result of $0/0$, as if it were a mean.

\begin{definition}[Defensive programming]\label{def:b1:c-modules-make-tests:defensive}
\emph{Defensive programming}\index{defensive programming} is writing each function so that it detects the
situations it was not designed for, instead of computing garbage from them: it checks what it receives from outside
the program and returns an error, and it asserts what its callers must guarantee.
\end{definition}

\begin{method}[Assertion or error return]\label{met:b1:c-modules-make-tests:assert-or-error}
\begin{enumerate}
\item If the situation can happen in a correct program (a file that does not exist, a line badly typed by a person, a
\code{malloc} that fails), return an error and let the caller decide: \code{date_parse} returns $-1$.
\item If it can happen only through a bug of the program (a month 13 passed to \code{days_in_month} after the date was
validated, a mean of no values), assert it: the program stops at the bug, with its location, instead of propagating a
wrong value.
\item Never put in an \omterm{def:b1:c-modules-make-tests:assertion}{assertion} an expression whose effect the program needs, such as \code{assert(p = malloc(n))}:
with \code{NDEBUG}, the call disappears.
\end{enumerate}
\end{method}

When a program crashes, the first question is where. Compiled with \code{-g}, the program carries the line numbers of
its source, and the address sanitizer prints the chain of calls that led to the faulty access, its backtrace.

\omcode{code/bachelor-1/07-c-modules-make-tests/listings/crash.c}{2}{27}{A search that does not stop at the end of the
list}\label{lst:b1:c-modules-make-tests:crash}

\omoutput{code/bachelor-1/07-c-modules-make-tests/outputs/backtrace.txt}{The crash, the access, and the two frames that
matter.}\label{out:b1:c-modules-make-tests:backtrace}

\begin{definition}[Debugger, breakpoint]\label{def:b1:c-modules-make-tests:debugger}
A \emph{debugger}\index{debugger} is a program that runs another program under its control: it can stop it, execute it
one line at a time, show and change its variables, and print the chain of active calls. A \emph{breakpoint}\index{breakpoint}
is a place in the program, usually a line, where the debugger stops it before executing that place.
\end{definition}

\begin{omfigure}
\begin{tikzpicture}[st/.style={draw, rounded corners=2pt, font=\scriptsize, align=left, inner sep=3pt, text width=31mm,
                    minimum height=15mm}, >={Stealth[length=2mm]}]
\node[st, fill=omCode!8] (a) at (0,0) {\textbf{1. breakpoint}\\on line 13 of \code{crash.c}, then start the
program};
\node[st, fill=omCode!8] (b) at (3.6,0) {\textbf{2. stop}\\before line 13; show \code{s->name}:
\code{"Hilltop"}};
\node[st, fill=omCode!8] (c) at (7.2,0) {\textbf{3. step}\\one line at a time; \code{s} becomes the address of
\code{r}, then \code{NULL}};
\node[st, fill=omThm!15] (d) at (10.8,0) {\textbf{4. backtrace}\\at the crash: \code{altitude_rank} line 13, called by
\code{main} line 25};
\draw[->] (a) -- (b); \draw[->] (b) -- (c); \draw[->] (c) -- (d);
\end{tikzpicture}
\par\smallskip{\small A debugging session on \cref{lst:b1:c-modules-make-tests:crash}, as a sequence of commands. The
program is built with \code{-g -O0}, so that each line of the source maps to machine instructions in order.}
\end{omfigure}

A typical session on the crash above, with a \omterm{def:b1:c-modules-make-tests:debugger}{debugger} such as the GNU \omterm{def:b1:c-modules-make-tests:debugger}{debugger}: set a \omterm{def:b1:c-modules-make-tests:debugger}{breakpoint} on line 13, start the
program, and it stops before the first comparison; print \code{s->name} and see \code{Hilltop}; step, and see \code{s}
move to Riverside, then to the \omterm{def:b1:pointers-and-arrays:null}{null pointer}; at the next step the program receives its signal, and the backtrace command
prints the two frames of the sanitizer's report. A \omterm{def:b1:c-modules-make-tests:debugger}{debugger} answers questions that a crash report cannot: what a
variable held three lines before, which branch was taken. No \omterm{def:b1:c-modules-make-tests:debugger}{debugger} is installed on the machine that ran this
book's programs, so this one is told, not shown; the sanitizer's backtrace above is available wherever the \omterm{def:b1:compilation-pipeline:compiler}{compiler} is.

\begin{safety}
Keep \omterm{def:b1:c-modules-make-tests:assertion}{assertions} on in tests and in development builds. Give every module its own test file, run the whole suite
before every change is kept, and add a test for each bug fixed. Build with \code{-Wall -Wextra -Werror}, with the address
sanitizer for tests, and with \code{-g} whenever a crash must be located.
\end{safety}

\section{Exercises}\label{sec:b1:c-modules-make-tests:exercises}

\begin{exercise}[$\star$]\label{exo:b1:c-modules-make-tests:1}
With the \omterm{def:b1:c-modules-make-tests:makefile}{makefile} of \cref{lst:b1:c-modules-make-tests:makefile}, after a full build, which commands does \code{make}
run after editing \code{check.h}? After editing \code{ratio.c}?
\end{exercise}

\begin{exercise}[$\star$]\label{exo:b1:c-modules-make-tests:2}
Does the test set $\{(4, 0), (70, 0)\}$ reach \omterm{def:b1:c-modules-make-tests:coverage}{statement coverage} of \code{fare}? Which single test makes it reach
\omterm{def:b1:c-modules-make-tests:coverage}{statement coverage}, and does the set then reach \omterm{def:b1:c-modules-make-tests:coverage}{branch coverage}?
\omcode{code/bachelor-1/07-c-modules-make-tests/listings/exo_fare.c}{2}{12}{A bus fare}
\end{exercise}

\begin{exercise}[$\star$]\label{exo:b1:c-modules-make-tests:3}
What does this program print, on which stream, and with which exit status? What if it is compiled with
\code{-DNDEBUG}?
\omcode{code/bachelor-1/07-c-modules-make-tests/listings/exo_assert.c}{2}{16}{A precondition}
\end{exercise}

\begin{exercise}[$\star$]\label{exo:b1:c-modules-make-tests:4}
\code{point.h} has no \omterm{def:b1:c-modules-make-tests:header}{include guard}. Explain the \omterm{def:b1:compilation-pipeline:compiler}{compiler}'s error, and add the guard.
\omcode{code/bachelor-1/07-c-modules-make-tests/listings/segment.h}{2}{6}{\code{segment.h} includes \code{point.h}}
\omcode{code/bachelor-1/07-c-modules-make-tests/listings/exo_guard.c}{2}{10}{So does the program}
\omoutput[2-3]{code/bachelor-1/07-c-modules-make-tests/outputs/exo_guard.txt}{The error.}
\end{exercise}

\begin{exercise}[$\star\star$]\label{exo:b1:c-modules-make-tests:5}
A colleague raised the limit in \code{config.h} from 3 to 5, ran \code{make}, and the program still says 3. Find the
missing line of the \omterm{def:b1:c-modules-make-tests:makefile}{makefile}.
\omcode{code/bachelor-1/07-c-modules-make-tests/listings/stale.Makefile}{2}{6}{A \omterm{def:b1:c-modules-make-tests:makefile}{makefile} with a missing
\omterm{def:b1:c-modules-make-tests:makefile}{prerequisite}}
\omoutput{code/bachelor-1/07-c-modules-make-tests/outputs/exo_stale.txt}{A stale build.}
\end{exercise}

\begin{exercise}[$\star\star$]\label{exo:b1:c-modules-make-tests:6}
The function \code{max3} has a bug, and its four tests pass. Find the bug, a test that reveals it, and what is wrong
with the four tests.
\omcode{code/bachelor-1/07-c-modules-make-tests/listings/exo_max3.c}{2}{21}{A green report}
\end{exercise}

\begin{exercise}[$\star\star$]\label{exo:b1:c-modules-make-tests:7}
Write a module \code{stack} of \code{int}: an opaque \code{struct stack}, \code{stack_new}, \code{stack_free},
\code{stack_push}, \code{stack_pop} (which reports an empty stack) and \code{stack_size}, with its header and a test
file that covers the empty stack and a stack that grows past its first capacity.
\end{exercise}

\begin{exercise}[$\star\star$]\label{exo:b1:c-modules-make-tests:8}
Give a set of tests that reaches \omterm{def:b1:c-modules-make-tests:coverage}{branch coverage} of \code{late_fee}, counting each atomic condition of \code{&&} as a
decision, and add the boundary of ``more than 30 days''.
\omcode{code/bachelor-1/07-c-modules-make-tests/exercises/late_fee.c}{4}{14}{A late fee, in cents}
\end{exercise}

\begin{exercise}[$\star\star$]\label{exo:b1:c-modules-make-tests:9}
A decision is \code{c1 && c2 && ... && ck}. How many tests try every combination of the $k$ conditions? Show
that $k + 1$ tests suffice for each condition to be shown, by a pair of tests differing only in that condition, to
change the outcome on its own. (This criterion, a form of modified condition/decision coverage, is told here, not
developed.)
\end{exercise}

\begin{exercise}[$\star\star\star$]\label{exo:b1:c-modules-make-tests:10}
In the \omterm{def:b1:c-modules-make-tests:makefile}{makefile} of the weekend problem, compute, for each of the seven \omterm{def:b1:compilation-pipeline:source-file}{source files}, the set of \omterm{def:b1:c-modules-make-tests:makefile}{targets} that
\code{make} rebuilds after an edit to that file, using only the graph of the rules. Explain why the answer is the set of
\omterm{def:b1:c-modules-make-tests:makefile}{targets} reachable from the edited file.
\end{exercise}

\begin{exercise}[$\star\star\star$]\label{exo:b1:c-modules-make-tests:11}
Write a \omterm{def:b1:c-modules-make-tests:makefile}{makefile} for the four modules of the weekend problem with a single pattern rule, a variable listing the
sources, and header dependencies written by the \omterm{def:b1:compilation-pipeline:compiler}{compiler} with \code{-MMD -MP}. Show that touching \code{csv.h} rebuilds
exactly \code{csv.o} and \code{loans.o}.
\end{exercise}

\begin{exercise}[$\star\star\star$]\label{exo:b1:c-modules-make-tests:12}
Suppose that every statement of a function lies on some \omterm{def:b1:c-modules-make-tests:cfg}{feasible path}, and that every decision has two outgoing arcs.
Prove that a test set reaching \omterm{def:b1:c-modules-make-tests:coverage}{branch coverage} reaches \omterm{def:b1:c-modules-make-tests:coverage}{statement coverage}. Give a function and a test set that reach
\omterm{def:b1:c-modules-make-tests:coverage}{statement coverage} but not \omterm{def:b1:c-modules-make-tests:coverage}{branch coverage}.
\end{exercise}

\section{Problem: Taming a one-file program}\label{sec:b1:c-modules-make-tests:problem}

\begin{problem}[{Weekend problem --- Taming a one-file program}]\label{pb:b1:c-modules-make-tests:1}
The library's loan checker has been split into four modules: \code{date} (calendar arithmetic), \code{csv} (splitting a
line), \code{loans} (a loan and its delay) and \code{loans_main.c} (the program). It reads the town's \code{loans.csv}
and a date, and reports the loans that were or are late at that date. The tests of the date module were written from
input classes; the date module's first version, \code{date_v0.c}, differs from the current one in one line of code.
\omcode{code/bachelor-1/07-c-modules-make-tests/listings/date.h}{2}{14}{The date module's \omterm{def:b1:c-modules-make-tests:interface}{interface}}\label{lst:b1:c-modules-make-tests:date-h}
\omcode{code/bachelor-1/07-c-modules-make-tests/listings/date.c}{6}{43}{Its \omterm{def:b1:c-modules-make-tests:interface}{implementation}}\label{lst:b1:c-modules-make-tests:date-c}
\omcode{code/bachelor-1/07-c-modules-make-tests/listings/loans.h}{2}{15}{The loans module's \omterm{def:b1:c-modules-make-tests:interface}{interface}}
\omcode{code/bachelor-1/07-c-modules-make-tests/listings/library.Makefile}{2}{24}{The \omterm{def:b1:c-modules-make-tests:makefile}{makefile}}\label{lst:b1:c-modules-make-tests:library-mk}
\omcode{code/bachelor-1/07-c-modules-make-tests/listings/date_test.c}{6}{36}{The tests of the date module}\label{lst:b1:c-modules-make-tests:date-test}
\omoutput{code/bachelor-1/07-c-modules-make-tests/outputs/problem.txt}{Building, running, counting the \omterm{def:b1:compilation-pipeline:compiler}{compiler} runs after
an edit to each file, testing, measuring coverage, and testing the first version.}\label{out:b1:c-modules-make-tests:problem}

\medskip\noindent\textbf{Part I --- Modules and interfaces.}
\begin{enumerate}
\item Draw the dependency graph of the checker: the seven \omterm{def:b1:compilation-pipeline:source-file}{source files}, the four \omterm{def:b1:compilation-pipeline:compiler}{object files} and the executable.
\item \code{loans.h} includes \code{date.h}. Why must it, and why does \code{loans.h} not include \code{csv.h}?
\item \code{days_since_1970} is \code{static} and absent from \code{date.h}. Give two reasons.
\item \code{struct date} is not opaque, unlike \code{struct ratio}. What would making it opaque cost the users of the
module, and what would it gain?
\item Which modules could be reused, unchanged, by a program that reads the weather file?
\item The comment of \code{csv_split} states its contract. Which part of it tells a caller that an empty last field
(a loan not returned) produces a field, and why does that matter to \code{loan_parse}?
\end{enumerate}

\medskip\noindent\textbf{Part II --- The makefile and the rebuilds it saves.}
\begin{enumerate}[resume]
\item The first build runs five commands. Which are \omterm{def:b1:compilation-pipeline:compiler}{compilations} and which is a link?
\item Explain each of the seven counts printed by the loop, from the graph of question 1.
\item Without \code{make}, a safe rebuild compiles the four \code{.c} files after any edit. Over one edit to each of the
seven files, how many \omterm{def:b1:compilation-pipeline:compiler}{compilations} is that, and how many did \code{make} perform?
\item The rule \code{loans_main.o: loans.h date.h} lists \code{date.h}, which \code{loans_main.c} does not include
itself. Is the line wrong?
\item Remove \code{date.h} from the rule of \code{loans.o}. Which edit would then produce a stale build, and what could
the stale program do wrong?
\item \code{date_test} is not part of the default \omterm{def:b1:c-modules-make-tests:makefile}{target}. What does that change in the counts of question 8?
\end{enumerate}

\medskip\noindent\textbf{Part III --- Tests from classes and boundaries.}
\begin{enumerate}[resume]
\item List the input classes of \code{is_leap} implied by the Gregorian rule, with one year of each from the tests.
\item Which tests of \code{date_parse} are \omterm{def:b1:c-modules-make-tests:input-class}{boundary values}? Name the boundary of each.
\item Recompute \code{between("1970-01-01", "2000-01-01") == 10957}.
\item The coverage run reaches every line and 30 of the 32 branches. Which two branches are never taken, and why can no
correct caller take them?
\item The checker reports 400 loans, 22 still out, 109 late and 5285 late days at 2025-12-31. Which of these numbers does
the \omterm{def:b1:c-modules-make-tests:unit-test}{test suite} check?
\item The date tests were written before the checker existed. Name one way in which they could still pass while the
checker gives a wrong total.
\end{enumerate}

\medskip\noindent\textbf{Part IV --- Defensive checks and a bug found by a test.}
\begin{enumerate}[resume]
\item The first version fails two checks. Which line of \code{date_v0.c} differs, and why do both failures come from it?
\item On the town's file, which covers 2025 only, does the first version give a different report? Why was the bug found
anyway?
\item \code{days_in_month} asserts that the month lies in 1 to 12, but \code{date_parse} returns $-1$ for month 13.
Explain this division of labour with the method of the chapter.
\item \code{loans_main.c} rejects a member number of 1000 or more. Is that a check or an \omterm{def:b1:c-modules-make-tests:assertion}{assertion}, and why?
\item What does the checker print, and with what exit status, if the date argument is \code{2025-02-29}?
\item State the named result of the problem, \emph{the rebuild saving}.
\end{enumerate}
\end{problem}
